\documentclass[10pt,a4paper]{amsart}
\usepackage[margin=19mm]{geometry}
\usepackage{amsmath,amssymb,mathtools}
\usepackage[scr=rsfs,cal=euler]{mathalfa}
\usepackage{xfrac}
\usepackage[unicode,hidelinks,bookmarks=false]{hyperref}
\newcommand{\ie}{i.e.\ }
\newcommand{\cf}{cf.\ }
\newcommand{\resp}{respectively\ }
\newcommand{\Subsection}{\subsection}
\providecommand{\nobreakdash}{\nobreak\hbox{-}}
\setlength{\emergencystretch}{2em}
\multlinegap0pt
\usepackage{bm}
\usepackage{bbm}
\usepackage{dsfont}
\usepackage{xspace}
\usepackage{cancel}
\usepackage{mathtools}
\usepackage[shortlabels]{enumitem}
\usepackage{amsthm}
\makeatletter
\@ifundefined{lemma}{\newtheorem{lemma}{Lemma}}{}
\@ifundefined{theorem}{\newtheorem{theorem}[lemma]{Theorem}}{}
\makeatother
\newcommand{\R}{{\mathcal{R}}}
\newcommand{\bC}{{\mathbb{C}}}
\newcommand{\bD}{{\mathbb{D}}}
\newcommand{\ol}{\overline}
\newcommand{\qand}{\quad\text{and}\quad}
\renewcommand{\phi}{\varphi}
\title[Annihilating ideals and Agler--McCarthy spectral varieties i]{Annihilating ideals and Agler--McCarthy spectral varieties in the bidisc}
\author{Rapha\"el Clou\^atre, Poornendu Kumar}
\date{Source: arXiv:2508.20826v2 (CC BY 4.0)}
\begin{document}
\maketitle

\noindent\emph{Source and licence.} Annihilating ideals and Agler--McCarthy spectral varieties in the bidisc. Version arXiv:2508.20826v2 (CC BY 4.0), distributed under the Creative Commons Attribution 4.0 International licence, as recorded in the arXiv licence metadata for this exact version and shown on the abstract page https://arxiv.org/abs/2508.20826v2. Full rights notice: licenses/arxiv-2508.20826-CC-BY-4.0.txt.\par\smallskip

\noindent\textit{Editorial scope.} Counted unit: the Theorem whose source label is \texttt{T:min2var} in the article (source0.tex lines 605--613, source label \texttt{T:min2var}), delivered together with the article's own complete original proof of it (source0.tex lines 614--642, one \texttt{proof} environment of 29 lines). No mathematics is written, repaired or re-derived: statement and proof are the article's own text line for line. Required context retained uncounted: T:Williams (lines 428--435); T:AMratspectral (lines 585--587). Every definition, display and label of those lines is retained unchanged. Omitted from this package and not redistributed: 1--427, 436--584, 588--604, 643--1143. Nothing inside a retained excerpt is omitted or altered: every retained body is delivered exactly as the article's lines are written, so no omission receipt of any kind accompanies this package. Citations: the retained excerpts carry no citation command at all, so no bibliography entry of the article is transcribed and no reference is redistributed. Reference adaptation: none was needed and none was made; every reference inside the delivered unit resolves inside it, so no unresolved reference or citation is produced. Printed slips kept literal: no printed word, symbol, hypothesis or number of the counted statement or of its proof was altered. Layout and class: the article's own class is \texttt{amsart} with packages including graphicx, amscd, amsmath, amsthm, amsfonts, amssymb, mathrsfs, bm, enumerate, amsrefs, xcolor, hyperref, inputenc, tikz; the wrapper is the lane's pinned \texttt{amsart} editorial wrapper at 10pt on A4, which restates the theorem-like environments the retained text needs and adds the article's own macro definitions that the retained text needs (\texttt{\textbackslash R}, \texttt{\textbackslash bC}, \texttt{\textbackslash bD}, \texttt{\textbackslash ol}, \texttt{\textbackslash phi}, \texttt{\textbackslash qand}), with extra wrapper packages (bm, bbm, dsfont, xspace, cancel, mathtools, enumitem). The delivered numbering is the unit's own. Assets: no retained span contains a figure environment, a graphics-inclusion command, a PDF-page inclusion, a table or a TikZ picture; no image file is copied into this package, the package declares no asset dependency and no image file is redistributed with it. Licence: version arXiv:2508.20826v2 of this article is distributed under the Creative Commons Attribution 4.0 International licence, as recorded in the arXiv licence metadata for this exact version and shown on the abstract page https://arxiv.org/abs/2508.20826v2. A full rights notice is delivered at licenses/arxiv-2508.20826-CC-BY-4.0.txt. Characters: every retained excerpt is pure ASCII, so the delivered engine, XeLaTeX with its default Latin Modern text font, reports no missing character.
\begin{theorem}\label{T:Williams} Let $G\subset \bC$ be a bounded, connected, open set with boundary consisting of finitely many rectifiable simple closed curves. Let $T\in B(H)$ be an operator admitting $\ol{G}$ as spectral set.
Assume that
\begin{enumerate}[{\rm (i)}]
\item $\sigma(T)$ is contained in $G$, and
\item there is a non-constant $\phi\in H^\infty(G)$ such that $\|\phi(T)\|=\|\phi\|_G=1$.
\end{enumerate}
If $\|\phi(T)\|\leq \|\phi\|_X$ for some proper closed subset $X\subset \ol{G}$, then $X$ is not a spectral set for $T$. In particular, $\ol{G}$ is a minimal spectral set for $T$.
\end{theorem}

\begin{theorem}\label{T:AMratspectral}
Let $T_1,T_2\in B(H)$ be commuting contractions. Assume that there is an Agler--McCarthy distinguished variety $V\subset \bD^2$ for $(T_1,T_2)$. Then, $\ol{V}$ is a spectral set for $(T_1,T_2)$.
\end{theorem}

\begin{theorem}\label{T:min2var}
Let $T_1,T_2\in B(H)$ be commuting contractions. Let $V\subset \bD^2$ be an Agler--McCarthy distinguished variety for $(T_1,T_2)$. 
Assume that
\begin{enumerate}[{\rm (i)}]
\item the spectrum of $T_1$ is contained in $\bD$, and
\item there are non-constant functions $\phi_1\in H^\infty(\bD)$ and $\phi_2\in \R(\ol{\bD})$ such that $1=\|\phi_1\|_{\bD}=\|\phi_2\|_{\bD}=\|\phi_1(T_1)\phi_2(T_2)\|$.
\end{enumerate}
Then $\ol{V}$ is a minimal spectral set for $(T_1,T_2)$.
\end{theorem}

\begin{proof}
It follows from Theorem \ref{T:AMratspectral} that $\ol{V}$ is a spectral set for $(T_1,T_2)$, so we need only show that it is minimal. There is a  rational matrix-valued pure inner function $\Psi$ on $\bD$ with the property that
\[
V=\{(z,w)\in \bD^2:\det(\Psi(z)-wI)=0\}.
\]
Assume that there is a proper closed subset $X\subset \ol{V}$ that is  a spectral set for $(T_1,T_2)$. 
Without loss of generality, we may assume that $X=\ol{V}\setminus D$ where $D\subset \bD^2$ is some polydisc of radius $\delta>0$ centred at some $(\alpha,\beta)\in V$. 


Next, $\phi_1$ can be approximated uniformly on compact subsets of $\bD$ by a sequence of polynomials $(p_n)$. Since the spectrum of $T_1$ is contained in $\bD$, it follows that $\phi_1(T_1)$ is the norm-limit of the sequence $(p_n(T_1))$, so that
\[
\|\phi_1(T_1)\|=\lim_{n\to \infty} \|p_n(T_1)\|\leq \lim_{n\to \infty} \|p_n\|_{\bD}=\|\phi_1\|_{\bD}=1
\]
as $\ol{\bD}$ is a spectral set for $T_1$ by von Neumann's inequality. Now, $\ol{\bD}$ is a spectral set for $T_2$ as well, so we find
\[
\|\phi_2(T)\|\leq \|\phi_2\|_{\bD}=1.
\]
The condition $\|\phi_1(T_1)\phi_2(T_2)\|=1$ then forces $\|\phi_1(T_1)\|=1=\|\phi_2(T_2)\|$.

Using that $X$ is spectral for $T$, there is $(\lambda,\mu)\in X$ such that
\[
1=\|\phi_1(T_1)\phi_2(T_2)\|\leq |\phi_1(\lambda)\phi_2(\mu)|.
\]
This implies 
\[
1=\|\phi_1(T_1)\|=|\phi_1(\lambda)| \qand 1=\|\phi_2(T_2)\|=|\phi_2(\mu)|.
\]
Because $X$ is a spectral set for $T$, it also follows that the set $X_1=\ol{\bD}\setminus B$ is a spectral set for $T_1$, where $B=\{z\in \bC:|z-\alpha|<\delta\}$.  By virtue of (i) we may apply Theorem \ref{T:Williams}  to see that  $\|\phi_1(T_1)\|>\|\phi\|_{X_1}$, so $\lambda\notin X_1$ and $\lambda\in B\subset \bD$. Using that $V$ is distinguished and $(\lambda, \mu)\in \ol{V}$ while $\lambda\in \bD$, we conclude that $\mu\in \bD$ as well,  so that $|\phi_2(\mu)|<1$ by the maximum modulus principle. This is a contradiction.
\end{proof}

\end{document}
